% ==========================================================================
% TEMPLATE DE RELATORIO FINAL - ESTAGIO EM DOCENCIA (PPGCC/DCC/UFMG)
% Estrutura exigida pelas Diretrizes (secao "Relatorio final"):
%   1. Disciplina/turma apoiada   2. Atividades conduzidas
%   3. Resultados obtidos         4. Autoavaliacao
%   5. Sugestoes para futuras ofertas
%
% COMO USAR:
%  1) Preencha os dados da capa nos comandos \newcommand logo abaixo.
%  2) Leia as caixas azuis "INSTRUCAO" em cada secao e substitua o texto
%     entre colchetes [ ... ] pelo seu conteudo.
%  3) Quando terminar, mude \mostrarinstrucoestrue para
%     \mostrarinstrucoesfalse para ocultar TODAS as caixas de instrucao.
%  4) Compile: pdflatex main; bibtex main; pdflatex main; pdflatex main
% ==========================================================================
\documentclass[12pt]{article}
\usepackage[brazilian]{babel}
\usepackage[utf8]{inputenc}
\usepackage[T1]{fontenc}
\usepackage{times}
\usepackage{amsmath, amssymb}
\usepackage{graphicx}
\usepackage{microtype}
\usepackage{xcolor}
\usepackage{booktabs}
\usepackage{float}
\usepackage{url}
\usepackage{tikz}
\usepackage{enumitem}
\usepackage{longtable}
\usepackage[most]{tcolorbox}
\usepackage[a4paper, top=2.5cm, bottom=2.5cm, left=2.5cm, right=2.5cm]{geometry}
\usepackage[colorlinks=true, linkcolor=dccblue, citecolor=dccblue, urlcolor=dccblue]{hyperref}

\definecolor{dccblue}{RGB}{20, 45, 110}
\definecolor{instrbg}{RGB}{232, 240, 254}

% ---------------------------- DADOS DA CAPA --------------------------------
\newcommand{\Semestre}{[AAAA/S]}
\newcommand{\NomeDisciplina}{[Nome da Disciplina apoiada]}
\newcommand{\CodigoDisciplina}{[CODIGO]}
\newcommand{\Turma}{[Turma]}
\newcommand{\NomeProfessor}{[Nome completo do Professor Supervisor]}
\newcommand{\EmailProfessor}{[email.professor@dcc.ufmg.br]}
\newcommand{\NomeEstagiario}{[Nome completo do Estagiario]}
\newcommand{\MatriculaEstagiario}{[Matricula UFMG]}
\newcommand{\NivelEstagiario}{[Mestrando/Doutorando] em Ciência da Computação}
\newcommand{\EmailEstagiario}{[email.estagiario@dcc.ufmg.br]}
\newcommand{\EstagioNivel}{[Estágio em Docência I ou II]}
\newcommand{\LocalData}{[Cidade, Mes de Ano]}

% ------------------ CAIXAS DE INSTRUCAO (ocultaveis) -----------------------
\newif\ifmostrarinstrucoes
\mostrarinstrucoestrue   % <- troque por \mostrarinstrucoesfalse na versao final
\newcommand{\instrucao}[1]{%
  \ifmostrarinstrucoes
  \begin{tcolorbox}[colback=instrbg, colframe=dccblue, arc=2mm, boxrule=0.6pt,
    title=\textbf{INSTRUÇÃO (apagar na versão final)}, fonttitle=\small]
  \small #1
  \end{tcolorbox}
  \fi}

\begin{document}

% ================================ CAPA ======================================
\begin{titlepage}
\newcommand{\HRule}{\rule{\linewidth}{0.5mm}}
\begin{tikzpicture}[remember picture, overlay]
  \node[anchor=south west, inner sep=1.2cm] at (current page.south west) {
    \includegraphics[height=1.25cm]{imgs/default/logo-ufmg.png}};
  \node[anchor=south east, inner sep=1.2cm] at (current page.south east) {
    \includegraphics[width=3.2cm]{imgs/default/Logo-DCC.png}};
\end{tikzpicture}
\center
\textsc{\Large Universidade Federal de Minas Gerais}\\[0.2cm]
\textsc{\large Instituto de Ciências Exatas --- ICEx}\\[0.2cm]
\textsc{\large Departamento de Ciência da Computação --- DCC}\\[0.2cm]
\textsc{\large Programa de Pós-Graduação em Ciência da Computação --- PPGCC}\\[0.35cm]
\includegraphics[scale=0.45]{imgs/default/logo.png}\\[0.35cm]
\textsc{\Large \EstagioNivel{} (\CodigoDisciplina)}\\[0.2cm]
\textsc{\large Semestre Letivo: \Semestre}\\[0.2cm]
\HRule \\[0.4cm]
{\huge \bfseries Relatório Final de Atividades de Estágio em Docência: \NomeDisciplina}\\[0.3cm]
\HRule \\[1.2cm]
\begin{minipage}{0.48\textwidth}
\begin{flushleft} \large
\emph{Professor Supervisor:}\\
\NomeProfessor\\
\small{\texttt{\EmailProfessor}}\\
\small{Departamento de Ciência da Computação}
\end{flushleft}
\end{minipage}
~
\begin{minipage}{0.48\textwidth}
\begin{flushright} \large
\emph{Estagiário em Docência:}\\
\NomeEstagiario\\
\small{Matrícula UFMG: \MatriculaEstagiario}\\
\small{\NivelEstagiario}\\
\small{\texttt{\EmailEstagiario}}
\end{flushright}
\end{minipage}\\[2.2cm]
\vfill
{\large \LocalData}
\end{titlepage}

\pagenumbering{roman}

% ================================ RESUMO ====================================
\begin{abstract}
\instrucao{Escreva um resumo de 1 parágrafo (150 a 250 palavras): disciplina e turma apoiada, papel do estagiário, principais atividades (elaboração/correção de trabalhos, atendimento a dúvidas, apoio em aulas e provas), principais resultados quantitativos e conclusão. Inclua 3 a 6 palavras-chave.}
\noindent [Resumo do relatório.]

\vspace{0.5cm}
\noindent \textbf{Palavras-chave:} [palavra 1], [palavra 2], [palavra 3].
\end{abstract}

\newpage
\tableofcontents
\newpage
\pagenumbering{arabic}

% ======================= 1. DISCIPLINA / TURMA APOIADA ======================
\section{Identificação e Caracterização da Disciplina Apoiada}\label{sec:disciplina}
\instrucao{Item 1 das Diretrizes: reproduza os dados da Proposta de Trabalho aprovada (identificação, condução da disciplina, ementa, bibliografia, avaliação e cronograma). O cronograma deve indicar o que foi planejado \emph{versus} o que foi efetivamente executado.}

\subsection{Identificação da Disciplina e da Turma}
\begin{itemize}[leftmargin=*]
    \item \textbf{Disciplina:} \NomeDisciplina; \textbf{Código:} \CodigoDisciplina; \textbf{Turma:} \Turma;
    \item \textbf{Horários:} [dias e horários das aulas];
    \item \textbf{Carga horária / créditos:} [ex.: 60 horas, 4 créditos];
    \item \textbf{Número de alunos matriculados:} [N];
    \item \textbf{Professor Supervisor:} \NomeProfessor;
    \item \textbf{Estagiário:} \NomeEstagiario.
\end{itemize}

\subsection{Condução Pedagógica pelo Professor Supervisor}
\instrucao{Descreva em 1 ou 2 parágrafos como o professor conduz a disciplina (aulas expositivas, laboratórios, trabalhos práticos, uso do Moodle etc.).}
[Descrição da condução da disciplina.]

\subsection{Ementa e Programa}
\instrucao{Copie a ementa oficial e liste os tópicos do programa da disciplina.}
\begin{quote}\textit{[Ementa oficial da disciplina.]}\end{quote}
\begin{enumerate}[noitemsep, topsep=2pt]
    \item [Tópico 1];
    \item [Tópico 2];
    \item [\ldots]
\end{enumerate}

\subsection{Bibliografia da Disciplina}
\instrucao{Liste a bibliografia básica e complementar adotada. Cadastre cada obra em \texttt{bibliografia.bib} e cite com \texttt{\textbackslash cite\{chave\}}.}
\begin{itemize}[leftmargin=*]
    \item \textbf{Básica:} [Autor (Título)] \cite{exemplo_livro};
    \item \textbf{Complementar:} [Autor (Título)].
\end{itemize}

\subsection{Sistema e Critérios de Avaliação}
\instrucao{Descreva a composição da nota da disciplina e as datas das avaliações. Adapte a tabela abaixo (linhas, valores e pesos).}
\begin{table}[H]
\centering\small
\caption{Composição das Formas de Avaliação}\label{tab:avaliacoes}
\vspace{0.2cm}
\begin{tabular}{llcc}
\toprule
\textbf{Avaliação} & \textbf{Descrição} & \textbf{Valor} & \textbf{Peso} \\
\midrule
{[AV1]} & [Prova / Trabalho --- data] & [00,00 pts] & [0,00] \\
{[AV2]} & [Prova / Trabalho --- data] & [00,00 pts] & [0,00] \\
\midrule
\textbf{Total} & & \textbf{[100,00 pts]} & --- \\
\bottomrule
\end{tabular}
\end{table}

\subsection{Cronograma Planejado versus Executado}
\instrucao{Inclua uma linha por aula/evento do semestre, indicando sua atuação. Marque feriados com ``---''. Informe se houve alteração em relação à Proposta de Trabalho.}
\begin{longtable}{|p{2.2cm}|p{7.0cm}|p{5.4cm}|}
\caption{Cronograma Semestral de Atividades de Estágio}\label{tab:cronograma}\\
\hline
\textbf{Data} & \textbf{Tópico / Evento} & \textbf{Atuação do Estagiário} \\ \hline
\endfirsthead
\hline
\textbf{Data} & \textbf{Tópico / Evento} & \textbf{Atuação do Estagiário} \\ \hline
\endhead
\hline
\endfoot
[DD/MM/AAAA] & [Aula teórica: tópico] & [Esclarecimento de dúvidas] \\ \hline
[DD/MM/AAAA] & [Feriado] & --- \\ \hline
[DD/MM/AAAA] & [Prova escrita] & [Apoio na aplicação da prova] \\ \hline
[DD/MM/AAAA] & [Encerramento do semestre] & [Entrega do Relatório Final] \\ \hline
\end{longtable}

\newpage
% ======================== 2. ATIVIDADES CONDUZIDAS ==========================
\section{Atividades Conduzidas pelo Estagiário}\label{sec:atividades}
\instrucao{Item 2 das Diretrizes: descreva, de forma objetiva e verificável, cada atividade sob sua responsabilidade (as mesmas previstas na Proposta de Trabalho). Sugestão de subseções abaixo; remova as que não se aplicam e acrescente outras.}

\subsection{Elaboração de Trabalhos e Avaliações}
\instrucao{Para cada trabalho/lista/prova elaborado: objetivo pedagógico, requisitos, formato de entrega e critérios de pontuação.}
\begin{itemize}[leftmargin=*]
    \item \textbf{[Trabalho 1 --- título]:} [descrição];
    \item \textbf{[Trabalho 2 --- título]:} [descrição].
\end{itemize}

\subsection{Correção e Atribuição de Notas}
\instrucao{Explique o método de correção (manual, automatizado, rubricas), prazos, ferramentas usadas, política de feedback e como as notas foram lançadas. Se houve verificação de integridade acadêmica, descreva a ferramenta e o critério adotado, \textbf{sem identificar alunos}.}
[Descrição do processo de correção.]

\subsection{Atendimento a Dúvidas e Monitoria}
\instrucao{Informe a modalidade (presencial, remota ou híbrida), canais utilizados, frequência e principais temas das dúvidas.}
[Descrição dos atendimentos.]

\subsection{Apoio em Aulas e Aplicação de Provas}
\instrucao{Liste as aulas práticas/laboratórios e provas em que você atuou, com datas e natureza do apoio.}
\begin{itemize}[leftmargin=*]
    \item [DD/MM/AAAA --- atividade e papel desempenhado].
\end{itemize}

\subsection{Outras Atividades}
\instrucao{Elaboração de materiais didáticos, ambientes de teste, comunicados aos alunos etc. Remova se não houver.}
[Outras atividades.]

\newpage
% ========================== 3. RESULTADOS OBTIDOS ===========================
\section{Resultados Obtidos}\label{sec:resultados}
\instrucao{Item 3 das Diretrizes: apresente evidências do resultado do seu trabalho (estatísticas da turma, evolução entre avaliações, taxa de entrega, desempenho médio, observações qualitativas). Use dados agregados --- \textbf{não inclua nomes nem matrículas de alunos}.}

\subsection{Desempenho Geral da Turma}
\begin{table}[H]
\centering\small
\caption{Quadro de Desempenho da Turma}\label{tab:desempenho}
\vspace{0.2cm}
\begin{tabular}{lcc}
\toprule
\textbf{Métrica} & \textbf{[Avaliação 1]} & \textbf{[Avaliação 2]} \\
\midrule
Alunos matriculados & [N] & [N] \\
Entregas / participações avaliadas & [N (\%)] & [N (\%)] \\
Média geral & [00,00] & [00,00] \\
Maior nota / menor nota & [00,00 / 00,00] & [00,00 / 00,00] \\
\bottomrule
\end{tabular}
\end{table}

\subsection{Análise e Discussão}
\instrucao{Interprete os números: dificuldades recorrentes, evolução ao longo do semestre, efeito das ações de apoio, ocorrência de recursos de nota.}
[Análise dos resultados.]

\newpage
% =========================== 4. AUTOAVALIAÇÃO ===============================
\section{Autoavaliação do Estagiário}\label{sec:autoavaliacao}
\instrucao{Item 4 das Diretrizes: reflexão crítica e honesta sobre sua atuação. Sugestão de eixos: competências docentes desenvolvidas, comunicação com os alunos, gestão de tempo/conciliação com a pós-graduação, pontos fortes, pontos a melhorar e contribuição para a vocação docente. O relatório é avaliado nos critérios: cumprimento dos objetivos, qualidade do apoio, disponibilidade e assiduidade, motivação e desenvoltura.}
\subsection{Competências Desenvolvidas}
[Texto.]
\subsection{Comunicação e Relacionamento com os Alunos}
[Texto.]
\subsection{Gestão de Tempo e Assiduidade}
[Texto.]
\subsection{Pontos a Melhorar e Reflexão Final}
[Texto.]

\newpage
% =================== 5. SUGESTÕES PARA FUTURAS OFERTAS ======================
\section{Sugestões para Futuras Ofertas da Disciplina}\label{sec:sugestoes}
\instrucao{Item 5 das Diretrizes: proponha melhorias concretas e justificadas. Para cada sugestão, indique o diagnóstico (o que foi observado) e a recomendação (o que fazer).}
\begin{itemize}[leftmargin=*]
    \item \textbf{[Sugestão 1]} --- \emph{Diagnóstico:} [observação]. \emph{Recomendação:} [ação];
    \item \textbf{[Sugestão 2]} --- \emph{Diagnóstico:} [observação]. \emph{Recomendação:} [ação].
\end{itemize}

\newpage
% ===================== CONSIDERAÇÕES FINAIS E ASSINATURAS ===================
\section{Considerações Finais e Termo de Encerramento}\label{sec:conclusao}
\instrucao{Declare que as atividades da Proposta de Trabalho foram cumpridas (ou justifique desvios), faça os agradecimentos e colete as assinaturas.}
[Considerações finais e agradecimentos.]

\vspace{2.5cm}
\noindent
\begin{minipage}[t]{0.48\textwidth}
    \centering
    \rule{\linewidth}{0.5pt}\\
    \textbf{\NomeProfessor}\\
    Professor Supervisor da Disciplina\\
    Departamento de Ciência da Computação --- DCC/UFMG
\end{minipage}
\hfill
\begin{minipage}[t]{0.48\textwidth}
    \centering
    \rule{\linewidth}{0.5pt}\\
    \textbf{\NomeEstagiario}\\
    \EstagioNivel\\
    PPGCC/UFMG
\end{minipage}

\newpage
\bibliographystyle{unsrt}
\bibliography{bibliografia}

\end{document}
